\pdfbookmark[2]{Zu VI und VII.}{nexus-remarks-vi-vii}
\begin{center}\large Zu VI und VII.\end{center}

Die Art, wie der Verf. die Summation der Reihe $\Sigma\left(\frac{-D}n\right)\frac1n$ ausgeführt hat, ergiebt sich aus einigen speciellen Beispielen, welche sich auf einzelnen Blättern vorfinden.

Ist $D\equiv3\pmod4$, so folgt aus dem Fundamentaltheorem in der Theorie der quadratischen Reste mit Benutzung der Reihe
\[
\operatorname{cotang}u=\frac1u+\frac1{u-\pi}+\frac1{u+\pi}+\frac1{u-2\pi}+\frac1{u+2\pi}+\ldots
\]
dass
\[
\Sigma\left(\frac{-D}n\right)\frac1n=\Sigma\left(\frac nD\right)\frac1n=\frac\pi{2D}\Sigma\left(\frac\nu D\right)\operatorname{cotang}\frac{\nu\pi}{2D}
\]
ist, wo $\nu$ alle relativen Primzahlen zu $2D$ durchläuft, die kleiner als $D$ sind; setzt man
\[
\surd-1=i,\quad\cos\frac{2\pi}D+i\sin\frac{2\pi}D=r
\]
und bezeichnet mit $\mu$ alle relativen Primzahlen zu $D$, welche nicht grösser als $D$ sind, so lässt sich die vorstehende Summe leicht in die folgende umformen
\[
\Sigma\left(\frac{-D}n\right)\frac1n=\frac{\pi i}{4D}\left(\frac2D\right)\Sigma\left(\frac\mu D\right)\frac{r^\mu-1}{r^\mu+1}
\]
Wendet man nun die für jede Wurzel $\omega$ der Gleichung $\omega^D=1$ gültige Formel
\SourceLeaf{308}
\[
\frac{\omega-1}{\omega+1}=\Sigma(-1)^{\alpha-1}\omega^\alpha
\]
an, in welcher $\alpha$ die Zahlen $1,2,3\ldots(D-1)$ durchlaufen muss, so erhält man durch Umkehrung der Summationsordnung
\[
\Sigma\left(\frac{-D}n\right)\frac1n=\frac{\pi i}{4D}\left(\frac2D\right)\Sigma(-1)^{\alpha-1}\Sigma\left(\frac\mu D\right)r^{\alpha\mu}
\]
Die auf $\mu$ bezügliche Summation lässt sich bekanntlich mit Hülfe der in der Abhandlung \emph{Summatio quarundam serierum singularium} bewiesenen Sätze ausführen; beschränkt man sich auf den Fall, in welchem $D$ durch kein Quadrat theilbar ist, so findet man allgemein
\[
\Sigma\left(\frac\mu D\right)r^{\alpha\mu}=\left(\frac\alpha D\right)i^{\left(\frac{D-1}2\right)^2}\surd D
\]
wo $\left(\frac\alpha D\right)=0$ gesetzt werden muss, falls $\alpha$ keine relative Primzahl zu $D$ ist. In dem Fall $D\equiv3\pmod4$ erhält man daher
\[
\Sigma\left(\frac{-D}n\right)\frac1n=\frac\pi{4\surd D}\left(\frac2D\right)\Sigma(-1)^\alpha\left(\frac\alpha D\right)=\frac\pi{2\surd D}\Sigma\left(\frac{\alpha^{\prime}}D\right)
\]
wo $\alpha^{\prime}$ alle relativen Primzahlen zu $D$ durchläuft, die kleiner als $\frac12D$ sind; da endlich $\epsilon=2$ ist, so wird die Anzahl der Classen
\[
k=\Sigma\left(\frac{\alpha^{\prime}}D\right)
\]

Ist dagegen $D\equiv1\pmod4$, so erhält man mit Benutzung der Reihe
\[
\operatorname{cosec}u=\frac1u-\frac1{u-\pi}-\frac1{u+\pi}+\frac1{u-2\pi}+\frac1{u+2\pi}-\ldots
\]
auf ähnliche Weise
\[
\begin{aligned}
\Sigma\left(\frac{-D}n\right)\frac1n&=\Sigma(-1)^{\frac{n-1}2}\left(\frac nD\right)\frac1n
=\frac\pi{2D}\Sigma(-1)^{\frac{\nu-1}2}\left(\frac\nu D\right)\operatorname{cosec}\frac{\nu\pi}{2D}\\
&=\frac\pi{2D}\Sigma\left(\frac\mu D\right)\frac{r^\mu}{r^{2\mu}+1}
\end{aligned}
\]
wo die Buchstaben $\nu$ und $\mu$ die frühere Bedeutung haben; schliesst man den evidenten Fall $D=1$ aus und wendet die für jede Wurzel $\omega$ der Gleichung $\omega^D=1$ (mit Ausnahme von $\omega=1$) gültige Formel
\[
\frac\omega{\omega\omega+1}=1+\Sigma\omega^{4\alpha^{\prime\prime}}+\Sigma\omega^{D-4\alpha^{\prime\prime}}
\]
an, in welcher $\alpha^{\prime\prime}$ die Zahlen $1,2,3\ldots\frac14(D-1)$ durchlaufen muss, so ergiebt sich, wieder unter der Beschränkung, dass $D$ durch kein Quadrat theilbar ist,
\[
\Sigma\left(\frac{-D}n\right)\frac1n=\frac\pi{\surd D}\Sigma\left(\frac{\alpha^{\prime\prime}}D\right)
\]
und hieraus, da $\epsilon=2$ ist,
\[
k=2\Sigma\left(\frac{\alpha^{\prime\prime}}D\right)
\]

\SourceLeaf{309}
Ganz ähnlich würden sich die Fälle behandeln lassen, in welchen $D$ gerade ist. —

Was die Bestimmung der Classen-Anzahl für \emph{positive} Determinanten $D$ betrifft, so finden sich ausser der im Text mitgetheilten Schlussformel nur einzelne geometrische Figuren vor, welche Hyperbel-Sectoren von endlichen Dimensionen darstellen, und neben denselben Ungleichungen, durch welche die Punkte, deren Coordinaten die Variabeln der quadratischen Formen sind, in das Innere eines solchen Hyperbel-Sectors gedrängt werden. Diese Hyperbel-Sectoren treten an die Stelle der Ellipsen, welche den quadratischen Formen von negativen Determinanten entsprechen, und durch die Bestimmung ihres Flächeninhalts ergiebt sich wieder die mittlere Darstellungsanzahl, wenn nämlich nur solche Darstellungen zugelassen werden, bei welchen die Variabeln den eben erwähnten Ungleichungen Genüge leisten. Andererseits dienen diese Ungleichungen dazu, aus den unendlich vielen Darstellungen einer Zahl $m$, welche alle zu einer und derselben Wurzel $n$ der Congruenz $nn-D\equiv0\pmod{\frac m{\mu\mu}}$ gehören und welche den sämmtlichen Auflösungen der Gleichung $tt-Duu=1$ entsprechen (vergl. Disqq. Arithm. art. 205), eine einzige zu isoliren und alle andern auszuschliessen. Die Anzahl aller zugelassenen Darstellungen der Zahl $m$ durch den Complex aller nicht eigentlich äquivalenten formae proprie primitivae ist dann gleich dem Werth der Function $f(m)$, in welcher nur $-D$ durch $D$ zu ersetzen ist, und aus der Betrachtung der Eigenschaften derselben ergiebt sich, wie früher bei negativen Determinanten, ein zweiter Ausdruck für die mittlere Darstellungsanzahl; die Vergleichung desselben mit dem vorher durch geometrische Betrachtungen abgeleiteten Werthe führt dann unmittelbar zu der Bestimmung der Anzahl der Classen.
\BookRule{35mm}
